% =============================================================================
% CAPÍTULO 6 — RESOLUCIÓN DE PROBLEMAS (TROUBLESHOOTING)
% =============================================================================
\chapter{Resolución de Problemas}
\label{ch:troubleshooting}
\resumenCapitulo{Este capítulo reúne los errores e incidencias más frecuentes que
pueden presentarse durante el uso de la plataforma, junto con sus acciones correctivas,
y responde a las preguntas más habituales de las personas usuarias.}

% -----------------------------------------------------------------------------
\section{Matriz de Errores Comunes y Acciones Correctivas}
\label{sec:matriz-errores}

La siguiente tabla resume los síntomas más frecuentes, su causa probable y la acción
recomendada para resolverlos.

\vspace{0.2cm}
\noindent
\renewcommand{\arraystretch}{1.4}
\fontsize{9}{11}\selectfont\sffamily
\begin{tabularx}{\textwidth}{|>{\bfseries\color{colorPrimario}}L{3.4cm}|Y|Y|}
\hline
\rowcolor{headerTabla}
\sffamily\textbf{Síntoma} & \sffamily\textbf{Causa probable} & \sffamily\textbf{Acción correctiva} \\
\hline
No es posible iniciar sesión & Credenciales incorrectas o cuenta sin verificar. &
Verifique el correo y la contraseña; use \boton{¿Olvidó su contraseña?} y confirme la
verificación enviada a su correo. \\
\hline
\rowcolor{grisTecnico}
No llega el código de verificación & El mensaje fue filtrado o el correo es incorrecto.
& Revise la carpeta de \textit{spam}, confirme la dirección y solicite un nuevo código. \\
\hline
La captura o el documento no se carga & Archivo demasiado grande o formato no admitido.
& Use formatos admitidos (PNG, JPG, PDF) y reduzca el tamaño del archivo. \\
\hline
\rowcolor{grisTecnico}
El mapa de ubicación no aparece & JavaScript deshabilitado o conexión inestable. &
Habilite JavaScript, recargue la página y verifique su conexión. \\
\hline
Los cambios no se guardan & No se pulsó el botón de guardado correspondiente. &
Asegúrese de pulsar \boton{Guardar} (o \boton{Actualizar}) tras editar cada sección. \\
\hline
\rowcolor{grisTecnico}
Un campo obligatorio impide continuar & Falta completar un campo marcado con asterisco
(*). & Complete todos los campos obligatorios antes de avanzar. \\
\hline
La sesión se cierra inesperadamente & Expiración por inactividad. & Vuelva a iniciar
sesión; active \campo{Recordarme} en su equipo personal. \\
\hline
\rowcolor{grisTecnico}
La interfaz se ve incompleta & Navegador obsoleto o caché corrupta. & Actualice el
navegador y limpie la caché; recargue con \comando{Ctrl+F5}. \\
\hline
\end{tabularx}

\vspace{0.3cm}
\begin{AvisoNota}
Ante cualquier comportamiento inesperado, el primer paso recomendado es recargar la
página y, si persiste, cerrar sesión y volver a iniciarla. Muchas incidencias
transitorias se resuelven así.
\end{AvisoNota}

% -----------------------------------------------------------------------------
\section{Preguntas Frecuentes (FAQ)}
\label{sec:faq}

\begin{description}[style=unboxed,leftmargin=0pt,noitemsep]
    \item[\textcolor{colorPrimario}{\textbf{¿Puedo cambiar mi tipo de perfil?}}]
    El cambio de rol (por ejemplo, de Profesional a Reclutador) lo realiza el equipo de
    administración. Contáctelo a través de los canales del capítulo~\ref{ch:soporte}.

    \item[\textcolor{colorPrimario}{\textbf{¿Mi portafolio es visible para todos?}}]
    Usted controla la visibilidad de su perfil y de cada proyecto mediante los
    interruptores de privacidad descritos en el capítulo~\ref{ch:profesional}.

    \item[\textcolor{colorPrimario}{\textbf{¿Cómo recupero datos eliminados?}}]
    La eliminación de proyectos, registros o cuentas es \textbf{permanente}. Antes de
    eliminar, exporte sus datos desde \rutanav{Configuración \textgreater\ Exportar mis datos}.

    \item[\textcolor{colorPrimario}{\textbf{¿Puedo usar EthosHub en el teléfono?}}]
    Sí. La plataforma cuenta con diseño adaptable; consulte el apartado de acceso móvil
    de cada capítulo de perfil.

    \item[\textcolor{colorPrimario}{\textbf{¿En qué idiomas está disponible?}}]
    La interfaz puede configurarse en español (ES), inglés (EN) y portugués (BR) desde
    \rutanav{Configuración \textgreater\ Personalización}.

    \item[\textcolor{colorPrimario}{\textbf{¿Cómo cambio entre tema claro y oscuro?}}]
    Desde \rutanav{Configuración \textgreater\ Personalización \textgreater\ Tema de la Interfaz},
    eligiendo \boton{Claro}, \boton{Oscuro} o \boton{Sistema}.
\end{description}
